% Provenance (never printed): NSI premiere "Algorithmique: algorithmes gloutons (rendu de monnaie, sac a dos,
% choix d'activites)". See sources/book2/BRIEFS.md.
\chapter{Greedy Algorithms}\label{ch:g11:greedy-algorithms}

A shopkeeper owes $87$ in change and has coins of $1$, $2$, $5$, $10$, $20$ and $50$. Without thinking she hands over
the largest coin that fits, $50$, then again the largest that fits what is left, $20$, then $10$, $5$, $2$: five
coins, and no way of giving fewer. Her rule never looks back and never looks ahead, and here it is the best possible.

Rules of this kind are called greedy. They are fast and easy to write, and they are often right, but not always: with
other coins the same shopkeeper gives one coin too many. This chapter writes greedy rules for three problems, proves
when they are right, and lets a program find where they go wrong.

\begin{recall}
\code{sorted(t, key=f)} sorts by the values of \code{f} (chapter \ref{ch:g10:data-tables}). The cost
of an algorithm counts its \omterm{def:g11:cost-of-an-algorithm:elementary}{elementary operations} (chapter \ref{ch:g11:cost-of-an-algorithm}); an invariant is kept by
each turn of a loop (chapter \ref{ch:g11:proving-algorithms-correct}).
\end{recall}

\section{Optimisation problems}\label{sec:g11:greedy-algorithms:optimisation}

\begin{definition}[Optimisation problem, objective, optimal solution]\label{def:g11:greedy-algorithms:optimisation}
An \emph{optimisation problem}\index{optimisation problem} asks, among the feasible solutions of a problem, for one
that makes a number as small or as large as possible. That number is the \emph{objective}\index{objective}. A
feasible solution with the best possible objective is an \emph{optimal solution}\index{optimal solution}.
\end{definition}

For change, a feasible solution is a set of coins that adds up to the amount, and the \omterm{def:g11:greedy-algorithms:optimisation}{objective} is the number of
coins, to be made as small as possible. On a small amount, a program can try every combination and keep the best.
That is slow, but surely right, and it serves as a reference.

\omcode{code/grade-11/07-greedy-algorithms/listings/change_exhaustive.py}{1}{25}{The fewest coins, by trying every
combination.}\label{lst:g11:greedy-algorithms:exhaustive}

\omoutput{code/grade-11/07-greedy-algorithms/outputs/change_exhaustive.txt}{$20$ needs $4$ coins of $1$, $2$ and $5$,
found among $29$ combinations; with six coins, greedy matches the fewest on every amount
checked.}\label{out:g11:greedy-algorithms:exhaustive}

\section{The greedy choice}\label{sec:g11:greedy-algorithms:greedy}

\begin{definition}[Greedy algorithm]\label{def:g11:greedy-algorithms:greedy}
A \emph{greedy algorithm}\index{greedy algorithm} builds a solution one piece at a time, each time taking the piece
that looks best at that moment, by a fixed rule, and never going back on a choice.
\end{definition}

\begin{method}[Designing a greedy algorithm]
\begin{enumerate}
  \item Choose the rule: the order in which candidates are looked at, often a sort by a well-chosen key.
  \item Take each candidate that keeps the solution feasible.
  \item Prove that the result is optimal, often by exchange: take any \omterm{def:g11:greedy-algorithms:optimisation}{optimal solution}, and change it, step by step,
    into the greedy one without making it worse.
  \item If there is no proof, look for a counter-example with a program that compares greedy with an exhaustive
    search on small inputs.
\end{enumerate}
\end{method}

\omcode{code/grade-11/07-greedy-algorithms/listings/change_greedy.py}{1}{15}{Greedy change: the largest coin that
fits, again and again.}\label{lst:g11:greedy-algorithms:change-greedy}

\omoutput{code/grade-11/07-greedy-algorithms/outputs/change_greedy.txt}{Three amounts with the coins $1$, $2$, $5$,
$10$, $20$, $50$.}\label{out:g11:greedy-algorithms:change-greedy}

With the coins $1$, $2$, $5$ an exchange argument proves greedy right. Take a way of paying with the fewest coins. It
uses at most one coin of $1$, since two could become one $2$; at most two coins of $2$, since three could become a $5$
and a $1$; and never two $2$ with a $1$, which could become a $5$. So the coins of $1$ and $2$ add up to at most $4$,
and every $5$ that fits is used, exactly as greedy does; what remains, $0$ to $4$, has only one way with these rules:
$0$, $1$, $2$, $2 + 1$, $2 + 2$, again greedy's way. Greedy and an \omterm{def:g11:greedy-algorithms:optimisation}{optimal solution} are the same.

For six coins a full case argument is long. A theorem does the work instead.

\begin{theorem}[Kozen and Zaks]\label{thm:g11:greedy-algorithms:kozen-zaks}
Let a system of coins have the values $1 = c_1 < c_2 < \dots < c_m$, with $m \ge 3$. If greedy change gives more
coins than necessary for some amount, then it does so already for some amount $x$ with $c_3 + 1 < x < c_m + c_{m-1}$.
\end{theorem}
\admitted
% ledger: g11:kozen-zaks

For $1$, $2$, $5$, $10$, $20$, $50$, it is enough to check every amount below $50 + 20 = 70$, which the output above
did: no failure there, so none anywhere. A system for which greedy change is always optimal has a name.

\begin{definition}[Canonical coin system]\label{def:g11:greedy-algorithms:canonical}
A coin system is \emph{canonical}\index{canonical coin system} if, for every amount, greedy change uses the fewest
coins possible.
\end{definition}

\section{Choosing activities}\label{sec:g11:greedy-algorithms:activities}

One room, twelve requests, each a start and an end hour; two bookings that overlap cannot both be kept. The goal is to
keep as many as possible. Several greedy rules come to mind: shortest first, earliest start first, earliest end first.
Only the last is always right.

\omcode{code/grade-11/07-greedy-algorithms/listings/activities.py}{1}{22}{Earliest end first: keep a request when it
starts after the last one kept ends.}\label{lst:g11:greedy-algorithms:activities}

\omoutput{code/grade-11/07-greedy-algorithms/outputs/activities.txt}{The requests sorted by end hour, and the four
kept.}\label{out:g11:greedy-algorithms:activities}

\begin{omfigure}
\begin{tikzpicture}[x=0.75cm, y=0.42cm, font=\footnotesize]
  \foreach \h in {8,...,21} {
    \draw[black!20] (\h,0.4) -- (\h,-12.4);
    \node[above] at (\h,0.4) {\h};
  }
  \foreach \n/\s/\e/\k [count=\r] in {G/10/11/1, I/10/11/0, D/11/12/1, B/15/17/1, H/15/17/0, A/15/18/0,
      E/16/18/0, C/16/19/0, J/16/19/0, K/18/19/1, L/17/19/0, F/18/21/0} {
    \ifnum\k=1 \def\fillc{omDef!45}\else\def\fillc{black!8}\fi
    \draw[fill=\fillc, draw=black!60] (\s,-\r+0.3) rectangle (\e,-\r-0.3);
    \node[left] at (8,-\r) {\n};
  }
\end{tikzpicture}

{\small The twelve requests, from the run, in order of end hour. In blue, the four that earliest end first keeps:
each starts after the previous blue one ends.}
\end{omfigure}

Why is it optimal? Let $g_1, g_2, \dots$ be the bookings greedy keeps, in order, and $o_1, o_2, \dots$ those of an
optimal choice, in order. Greedy ``stays ahead'': $g_1$ ends no later than $o_1$, since greedy took the request that
ends first. If $g_k$ ends no later than $o_k$, then $o_{k+1}$, which starts after $o_k$ ends, starts after $g_k$ ends
too; so it was still available to greedy, which took a request ending no later, and $g_{k+1}$ ends no later than
$o_{k+1}$. Greedy can therefore never run out before the optimal choice does: it keeps at least as many bookings.

\begin{omfigure}
\begin{tikzpicture}[x=0.8cm, y=0.6cm, font=\footnotesize]
  \node[left] at (0,0) {optimal}; \node[left] at (0,-1.4) {greedy};
  \foreach \s/\e/\l in {0.5/2.5/{$o_1$}, 3/5.5/{$o_2$}, 6.5/8.5/{$o_3$}}
    \draw[fill=black!8] (\s,0.3) rectangle node {\l} (\e,-0.3);
  \foreach \s/\e/\l in {0.2/2/{$g_1$}, 2.5/4.6/{$g_2$}, 5/7/{$g_3$}}
    \draw[fill=omDef!40] (\s,-1.1) rectangle node {\l} (\e,-1.7);
  \draw[dashed, omThm] (2,-2) -- (2,0.7); \draw[dashed, omThm] (2.5,-2) -- (2.5,0.7);
  \draw[->, thick, omThm] (2,-2.3) -- node[below, font=\scriptsize] {ends no later} (2.5,-2.3);
\end{tikzpicture}

{\small Greedy stays ahead. Each $g_k$ ends no later than $o_k$, so $o_{k+1}$ was still free for greedy, and
replacing $o_k$ by $g_k$ in the optimal choice keeps it feasible.}
\end{omfigure}

\section{Filling a knapsack}\label{sec:g11:greedy-algorithms:knapsack}

A trader's sack holds $15$~kg. Goods can be cut: half a sack of rice is worth half its price. The greedy rule takes the
goods in order of value per kilogram, as much as fits.

\omcode{code/grade-11/07-greedy-algorithms/listings/knapsack_fractional.py}{1}{21}{The fractional knapsack: the best
value per kilogram first.}\label{lst:g11:greedy-algorithms:fractional}

\omoutput{code/grade-11/07-greedy-algorithms/outputs/knapsack_fractional.txt}{Three goods whole, then $3$~kg of rice
fill the sack.}\label{out:g11:greedy-algorithms:fractional}

\begin{omfigure}
\begin{tikzpicture}[x=0.4cm, y=0.12cm, font=\footnotesize]
  \draw[->] (0,0) -- (31,0) node[below left] {kg};
  \draw[->] (0,0) -- (0,33) node[above] {value per kg};
  \foreach \x/\w/\d/\n in {0/2/25/spices, 2/4/15/tea, 6/6/12/coffee, 12/10/4/rice, 22/8/2/flour} {
    \pgfmathsetmacro{\e}{min(\x+\w, 15)}
    \draw[fill=black!8] (\x,0) rectangle (\x+\w,\d);
    \ifdim \x pt<15pt \draw[fill=omDef!45] (\x,0) rectangle (\e,\d); \fi
    \node[above, font=\scriptsize] at (\x+\w/2,\d) {\n};
  }
  \draw[thick, omThm] (15,0) -- (15,26) node[above, font=\scriptsize] {sack full: 15 kg};
\end{tikzpicture}

{\small The goods as bars, widths in kilograms, heights in value per kilogram, best first. The sack takes everything
left of the line; the value is the blue area, $194$.}
\end{omfigure}

Here greedy is optimal: any kilogram in the sack that is not among the best ones could be swapped for a better one
left out, which raises the value. Now suppose the objects cannot be cut: a lamp is taken whole or not at all.

\omcode{code/grade-11/07-greedy-algorithms/listings/knapsack_01_greedy_fails.py}{1}{18}{Whole objects: the same rule,
and every subset as a reference.}\label{lst:g11:greedy-algorithms:knapsack-01}

\omoutput{code/grade-11/07-greedy-algorithms/outputs/knapsack_01_greedy_fails.txt}{The densest object takes $6$~kg of
the $10$ and leaves room for nothing else.}\label{out:g11:greedy-algorithms:knapsack-01}

No greedy rule is known to solve the knapsack with whole objects; the grade-12 chapter on dynamic programming solves it
exactly by a different method.

\section{When greedy fails}\label{sec:g11:greedy-algorithms:fails}

With coins of $1$, $3$ and $4$, greedy pays $6$ as $4 + 1 + 1$, three coins, while $3 + 3$ needs two. The program below
looks, among all systems of three coins up to $8$, for those that are not \omterm{def:g11:greedy-algorithms:canonical}{canonical}, and for the first amount each
gets wrong; the theorem says to look only below the sum of the two largest coins.

\omcode{code/grade-11/07-greedy-algorithms/listings/smallest_counterexample.py}{1}{13}{Every system of three coins up to
$8$ where greedy fails.}\label{lst:g11:greedy-algorithms:smallest}

\omoutput{code/grade-11/07-greedy-algorithms/outputs/smallest_counterexample.txt}{The smallest system that fails is
$1$, $3$, $4$, first at $6$.}\label{out:g11:greedy-algorithms:smallest}

\begin{omfigure}
\begin{tikzpicture}[coin/.style={circle, draw, thick, minimum size=9mm, font=\small}]
  \node[left] at (-0.4,0) {greedy};
  \node[coin, fill=omThm!20] at (0.5,0) {4};
  \node[coin, fill=omMeth!20] at (1.6,0) {1};
  \node[coin, fill=omMeth!20] at (2.7,0) {1};
  \node[right] at (3.4,0) {$3$ coins};
  \node[left] at (-0.4,-1.3) {fewest};
  \node[coin, fill=omDef!20] at (0.5,-1.3) {3};
  \node[coin, fill=omDef!20] at (1.6,-1.3) {3};
  \node[right] at (3.4,-1.3) {$2$ coins};
\end{tikzpicture}

{\small The amount $6$ with coins $1$, $3$, $4$. Taking $4$ first leaves $2$, which only coins of $1$ can pay.}
\end{omfigure}

\section{Exercises}\label{sec:g11:greedy-algorithms:exercises}

\begin{exercise}[$\star$]\label{exo:g11:greedy-algorithms:1}
Give greedy change for $87$ with coins $1$, $2$, $5$, $10$, $20$, $50$, without the program.
\end{exercise}

\begin{exercise}[$\star$]\label{exo:g11:greedy-algorithms:2}
With coins $1$, $3$, $4$, how many coins does greedy give for $6$? How many are needed?
\end{exercise}

\begin{exercise}[$\star$]\label{exo:g11:greedy-algorithms:3}
Four requests: P from $8$h to $12$h, S from $11$h to $13$h, Q from $12$h to $16$h, R from $13$h to $14$h. Which does
earliest end first keep?
\end{exercise}

\begin{exercise}[$\star$]\label{exo:g11:greedy-algorithms:4}
With the goods of the chapter and a sack of $8$~kg, what does the fractional greedy take, and for what value?
\end{exercise}

\begin{exercise}[$\star$]\label{exo:g11:greedy-algorithms:5}
Write \code{greedy_change(amount, coins)} with \code{divmod}, one division per coin value.
\end{exercise}

\begin{exercise}[$\star\star$]\label{exo:g11:greedy-algorithms:6}
A student sorts the requests by start time instead of end time. Find four requests on which this keeps one booking
when three are possible.
\end{exercise}

\begin{exercise}[$\star\star$]\label{exo:g11:greedy-algorithms:7}
Write \code{max_activities(requests)}, each request being a triple \code{(name, start, end)}, that returns the names
kept by earliest end first.
\end{exercise}

\begin{exercise}[$\star\star$]\label{exo:g11:greedy-algorithms:8}
How many combinations does \code{fewest(20, [1, 2, 5])} try? Count them by hand, by the number of coins of $5$, then of
$2$.
\end{exercise}

\begin{exercise}[$\star\star$]\label{exo:g11:greedy-algorithms:9}
With a sack of $10$~kg, compare the fractional greedy by value per kilogram with a greedy that takes the most valuable
goods first. Which is better here, and why?
\end{exercise}

\begin{exercise}[$\star\star$]\label{exo:g11:greedy-algorithms:10}
Objects that cannot be cut: X of $5$~kg worth $40$, Y and Z of $4$~kg worth $28$ each; the sack holds $8$~kg. What does
greedy by value per kilogram keep, and what is optimal?
\end{exercise}

\begin{exercise}[$\star\star$]\label{exo:g11:greedy-algorithms:11}
Write \code{is_canonical_upto(coins, m)} that compares greedy with \code{fewest} on every amount from $1$ to $m$. With
the theorem, which $m$ makes it a full test?
\end{exercise}

\begin{exercise}[$\star\star\star$]\label{exo:g11:greedy-algorithms:12}
Write out in full the proof that earliest end first keeps as many bookings as possible, with the ``stays ahead''
statement as an invariant of its loop.
\end{exercise}

\begin{exercise}[$\star\star\star$]\label{exo:g11:greedy-algorithms:13}
Prove that greedy change is optimal for the coins $1$, $2$, $5$ by treating the amounts left after the coins of $5$:
$0$, $1$, $2$, $3$ and $4$.
\end{exercise}

\begin{exercise}[$\star\star\star$]\label{exo:g11:greedy-algorithms:14}
All the requests must now be kept, in as few rooms as possible. Taking the requests by start time, give each the first
room free at its start, opening a new room if none is. Write \code{rooms_needed(requests)} and check that it equals the
largest number of requests going on at the same hour.
\end{exercise}

\begin{exercise}[$\star\star\star$]\label{exo:g11:greedy-algorithms:15}
Find three requests on which ``shortest first'' keeps fewer bookings than possible.
\end{exercise}

\section{Problem: The Coins of Brenmoor}\label{sec:g11:greedy-algorithms:problem}

\begin{problem}[{Weekend problem --- the amounts greedy gets wrong}]\label{pb:g11:greedy-algorithms:1}
The town of Brenmoor, invented for this problem, mints coins of $1$, $2$, $5$, $12$, $25$ and $40$. Its shopkeepers give
change greedily. The program compares greedy with the exhaustive search on every amount up to $100$, then tries every
way of changing one coin.

\omcode{code/grade-11/07-greedy-algorithms/listings/brenmoor.py}{1}{24}{Greedy against the fewest coins in Brenmoor, and
the systems one coin away.}\label{lst:g11:greedy-algorithms:brenmoor}

\omoutput{code/grade-11/07-greedy-algorithms/outputs/brenmoor.txt}{The failures, and four ways to remove them all by
changing one coin.}\label{out:g11:greedy-algorithms:brenmoor}

\needspace{4\baselineskip}\medskip\noindent\textbf{Part I --- The coins.}
\begin{enumerate}
  \item Give greedy change for $24$, $30$ and $37$.
  \item Why must a coin system contain a coin of $1$ for greedy to pay every amount?
  \item By the theorem, below which amount must a failure appear, if there is one?
  \item Why does the program still check up to $100$?
\end{enumerate}

\needspace{4\baselineskip}\medskip\noindent\textbf{Part II --- Greedy against the fewest.}
\begin{enumerate}[resume]
  \item Give greedy change for $36$, and a way with fewer coins.
  \item By the theorem's lower bound, which amounts cannot be the first failure? Which amounts did the program have to
    examine before finding $36$?
  \item For $49$ and $50$, give greedy's change and a better one.
  \item The exhaustive search tries $3687$ combinations for $100$. Why does greedy need only a few steps?
  \item Is the system \omterm{def:g11:greedy-algorithms:canonical}{canonical}?
\end{enumerate}

\needspace{4\baselineskip}\medskip\noindent\textbf{Part III --- The failures.}
\begin{enumerate}[resume]
  \item What do $36$, $61$ and $76$ have in common in their best change?
  \item For $60$, $61$, $62$ and $63$, give greedy's change and a way with fewer coins.
  \item How many coins too many does greedy give for $100$? Find a way with only four.
  \item How many of the amounts $1$ to $100$ fail?
\end{enumerate}

\needspace{4\baselineskip}\medskip\noindent\textbf{Part IV --- Changing one coin.}
\begin{enumerate}[resume]
  \item Replacing $12$ by $10$ gives a \omterm{def:g11:greedy-algorithms:canonical}{canonical} system. Up to which amount did the program have to check, and why is
    that enough?
  \item Check by hand that $36$ no longer fails with $1$, $2$, $5$, $10$, $25$, $40$.
  \item Replacing $25$ by $26$ also works. Check $36$ and $50$ in that system.
  \item Why does the program skip replacements that give a coin already in the system?
  \item Brenmoor prefers to keep its $12$ and its $40$. Which replacements allow it?
  \item Name an advantage for the town of a \omterm{def:g11:greedy-algorithms:canonical}{canonical} system.
  \item State the result: the smallest amount greedy gets wrong in Brenmoor, and how many amounts up to $100$ it gets
    wrong.
\end{enumerate}
\end{problem}
